% ==============================================================================
% CHƯƠNG 3: KIẾN TRÚC SUY DIỄN NEURO-SYMBOLIC, ĐẠI SỐ HÓA VÀ DỰNG HÌNH PHỤ
% ==============================================================================

\chapter{KIẾN TRÚC SUY DIỄN NEURO-SYMBOLIC, ĐẠI SỐ HÓA VÀ DỰNG HÌNH PHỤ}
\label{chap:reasoning}

\section{Động cơ Suy diễn Tiến Ký hiệu (Forward Chaining Deduction Engine)}
Cốt lõi của cỗ máy chứng minh hình thức là \textbf{Động cơ Suy diễn Tiến (Forward Chaining Deduction Engine)}, hoạt động trên Bộ nhớ Làm việc ($\mathcal{WM}$ -- Working Memory).

\subsection{Cấu trúc Tập Luật và Hợp giải Vị từ (Unification)}
Mỗi luật suy diễn hình học $r \in \mathcal{R}$ được định nghĩa bởi cặp tiền đề và kết luận:
\begin{equation}
    r = \langle \text{Inputs}(r), \text{Outputs}(r) \rangle
\end{equation}
trong đó $\text{Inputs}(r) = \{p_1(?x_1, \dots), \dots, p_k(?x_k, \dots)\}$ là tập các mẫu vị từ chứa biến (bắt đầu bằng dấu \texttt{?}), và $\text{Outputs}(r) = \{q_1(?x_1, \dots), \dots, q_m(?x_m, \dots)\}$ là các sự kiện suy ra sau khi thế biến.

\subsection{Chuẩn hóa Đồng nhất và Quan hệ Tương đương (Canonical Equivalence)}
Trong hình học, một đối tượng có thể được biểu diễn qua nhiều hoán vị đỉnh có tính chất tương đương (ví dụ: đoạn thẳng $AB \equiv BA$, góc $\angle ABC \equiv \angle CBA$, tam giác $\triangle ABC \equiv \triangle BCA$). Nếu không chuẩn hóa, bộ nhớ làm việc sẽ bị bùng nổ tổ hợp do lưu trữ trùng lặp.

Hệ thống định nghĩa hàm chuẩn hóa chính quy $\text{Canonical}(P)$ cho mọi vị từ:
\begin{itemize}
    \item Với đoạn thẳng: $\text{Canonical}(\text{Segment}(A, B)) = \text{Segment}(\min(A, B), \max(A, B))$.
    \item Với góc: $\text{Canonical}(\text{Angle}(A, B, C)) = \text{Angle}(\min(A, C), B, \max(A, C))$.
    \item Với đường thẳng song song:
    $$\text{Parallel}(AB, CD) \equiv \text{Parallel}(CD, AB) \equiv \text{Parallel}(BA, CD) \equiv \text{Parallel}(DC, AB)$$
\end{itemize}

\begin{algorithm}[h!]
\caption{Thuật toán Suy diễn Tiến Hình học Phẳng (Forward Chaining with Canonical Normalization)}
\label{alg:forward_chaining}
\begin{algorithmic}[1]
\Require Tập sự kiện ban đầu $\mathcal{F}_0$, Tập luật $\mathcal{R}$, Vị từ mục tiêu $\mathcal{G}$, Ngưỡng lặp cực đại $M$
\Ensure Trạng thái chứng minh (\texttt{True/False}), Dấu vết suy diễn $\mathcal{T}$
\State $\mathcal{WM} \leftarrow \bigcup_{f \in \mathcal{F}_0} \text{Canonical}(f)$
\State $\mathcal{T} \leftarrow []$
\State $step \leftarrow 0$
\While{$step < M$}
    \If{$\text{CheckGoal}(\mathcal{WM}, \mathcal{G})$}
        \State \Return \textbf{True}, $\mathcal{T}$ \Comment{Đã đạt được mục tiêu chứng minh}
    \EndIf
    \State $new\_facts \leftarrow \emptyset$
    \For{mỗi luật $r \in \mathcal{R}$}
        \State $bindings \leftarrow \text{UnifyMultiple}(\text{Inputs}(r), \mathcal{WM})$
        \For{mỗi phép thế $\theta \in bindings$}
            \State $conclusions \leftarrow \text{Substitute}(\text{Outputs}(r), \theta)$
            \For{mỗi $c \in conclusions$}
                \State $c_{can} \leftarrow \text{Canonical}(c)$
                \If{$c_{can} \notin \mathcal{WM}$}
                    \State $new\_facts \leftarrow new\_facts \cup \{c_{can}\}$
                    \State $\mathcal{T} \leftarrow \mathcal{T} \cup \{\langle r, \theta, c_{can} \rangle\}$
                \EndIf
            \EndFor
        \EndFor
    \EndFor
    \If{$new\_facts = \emptyset$}
        \State \textbf{break} \Comment{Đạt trạng thái bão hòa (Bế tắc)}
    \EndIf
    \State $\mathcal{WM} \leftarrow \mathcal{WM} \cup new\_facts$
    \State $step \leftarrow step + 1$
\EndWhile
\State \Return \textbf{False}, $\mathcal{T}$
\end{algorithmic}
\end{algorithm}

\section{Động cơ Suy diễn Đại số SymPy AR (Algebraic Reasoning)}
Các quan hệ về góc (Angle Chasing) và độ dài đoạn thẳng thường đòi hỏi các phép biến đổi phương trình liên tục mà suy diễn vị từ rời rạc khó có thể bao quát hết mọi phép biến đổi tương đương. Do đó, hệ thống tích hợp động cơ đại số hình thức \textbf{SymPy AR} chạy song song với động cơ suy diễn tiến.

\subsection{Cơ chế Hoạt động của SymPy AR}
\begin{enumerate}
    \item \textbf{Đại số hóa Vị từ (Equation Extraction)}: Chuyển các vị từ cộng góc $\text{Equal}(\text{Add}(\text{Angle}(A), \text{Angle}(B)), \text{Angle}(C))$ thành các phương trình đại số dạng biểu thức tượng trưng SymPy:
    $$\text{Equal}(\text{Add}(\text{Angle}(A), \text{Angle}(B), \text{Angle}(C)), 180) \implies \theta_A + \theta_B + \theta_C - 180 = 0$$
    \item \textbf{Khử Ẩn số và Giải Hệ Phương trình (Elimination \& Linear Solving)}: Sử dụng ma trận giải hệ phương trình tuyến tính hoặc phương pháp phần bù (Elimination Method) của SymPy để giải tìm giá trị ẩn số của mục tiêu:
    $$\theta_C = 180 - \theta_A - \theta_B = 180 - 60 - 70 = 50^\circ$$
    \item \textbf{Tái tạo Vị từ (Fact Injection)}: Khi một giá trị đại số mới được chứng minh, nó được đóng gói lại thành vị từ chính quy $\text{Equal}(\text{Angle}(C), 50)$ và nạp ngược trở lại $\mathcal{WM}$.
\end{enumerate}

\section{Tác tử Sinh Hình phụ Neuro-Symbolic (Auxiliary Construction Agent)}
Trong các bài toán hình học nâng cao (đặc biệt là đề thi học sinh giỏi và Olympic IMO), nhiều định lý không thể chứng minh trực tiếp nếu chỉ dựa trên các đối tượng ban đầu. Cần phải kẻ thêm các \textbf{đối tượng hình học phụ (Auxiliary Elements)}.

\begin{figure}[h!]
\centering
\begin{tikzpicture}[scale=0.88, transform shape]
    % Left column
    \node [base, fill=blue!12, draw=blue!85!black, minimum width=3.8cm, minimum height=1.3cm, inner sep=6pt] (start) at (-4.5, 2.2) {
        \textbf{Tập sự kiện ban đầu $\mathcal{F}_0$}\\[3pt]
        {\scriptsize (Working Memory $\mathcal{WM}$)}
    };
    
    \node [agent, fill=purple!15, draw=purple!85!black, minimum width=3.8cm, minimum height=1.3cm, inner sep=6pt] (aux) at (-4.5, -2.4) {
        \textbf{Auxiliary Agent (LLM)}\\[3pt]
        {\scriptsize Phân tích Gap \& Dựng hình phụ}
    };

    % Center column
    \node [process, fill=green!15, draw=green!65!black, minimum width=4.5cm, minimum height=1.3cm, inner sep=6pt] (forward) at (1.0, 2.2) {
        \textbf{Forward Deduction Engine}\\[3pt]
        {\scriptsize (Luật vị từ + FormalGeo)}
    };
    
    \node [process, fill=yellow!18, draw=orange!75!black, minimum width=4.5cm, minimum height=1.3cm, inner sep=6pt] (sympy) at (1.0, -0.1) {
        \textbf{SymPy AR Engine}\\[3pt]
        {\scriptsize (Đại số hóa góc \& đoạn thẳng)}
    };
    
    \node [decision, fill=yellow!18, draw=orange!85!black, minimum width=3.0cm, minimum height=1.4cm] (goal_check) at (1.0, -2.4) {
        \textbf{Đạt}\\[1pt]\textbf{Goal?}
    };

    % Right column
    \node [base, fill=teal!18, draw=teal!85!black, minimum width=3.6cm, minimum height=1.3cm, inner sep=6pt] (done) at (6.2, -2.4) {
        \textbf{Hoàn thành! \faCheckCircle}\\[3pt]
        {\scriptsize Xuất lời giải logic}
    };

    % Arrows with clear, non-overlapping labels
    \draw [arrow] (start) -- (forward);
    \draw [arrow] (forward) -- (sympy);
    \draw [arrow] (sympy) -- (goal_check);
    \draw [arrow] (goal_check.east) -- node[above=3pt, font=\bfseries\small, text=teal!85!black] {Có} (done.west);
    \draw [arrow] (goal_check.west) -- node[above=3pt, font=\bfseries\small, text=red!85!black] {Bế tắc} (aux.east);
    \draw [arrow] (aux.north) -- node[left=4pt, font=\itshape\footnotesize, text=purple!90!black, align=right] {Nạp hình phụ\\vào $\mathcal{WM}$} (start.south);
\end{tikzpicture}
\caption{Quy trình phối hợp toàn diện giữa Suy diễn tiến, SymPy AR và Tác tử Sinh hình phụ.}
\label{fig:full_reasoning_flow}
\end{figure}

\subsection{Chiến lược Đề xuất Hình phụ}
Khi cỗ máy suy diễn đạt điểm bất động mà chưa đạt mục tiêu, \textbf{Tác tử Bổ trợ Hình phụ (Auxiliary Agent)} được kích hoạt:
\begin{enumerate}
    \item \textbf{Nhận diện Khoảng trống Chứng minh (Gap Analysis)}: So sánh tập sự kiện hiện có với cấu trúc vị từ của mục tiêu (ví dụ: mục tiêu yêu cầu chứng minh $\text{CyclicQuadrilateral}$, nhưng hiện chưa có đường tròn ngoại tiếp).
    \item \textbf{Đề xuất Mẫu Hình phụ Hợp lệ}:
    \begin{itemize}
        \item Dựng đường kính hoặc tâm đường tròn ngoại tiếp: $\text{Circumcircle}(\text{Circle}(O), \text{Triangle}(A, B, C))$.
        \item Lấy trung điểm đoạn thẳng: $\text{Midpoint}(M, AB)$.
        \item Dựng đường cao hoặc đường vuông góc: $\text{Perpendicular}(AH, BC)$.
        \item Nối hai đỉnh chưa liên kết để tạo tam giác đồng dạng.
    \end{itemize}
\end{enumerate}
